\section{Detalle completo de la revisión}
\label{app:review-full-detail}

Las Secciones~\ref{subsec:method-review-protocol} y~\ref{subsec:tf-discriminant}
presentan las tres revisiones en formato mínimo: una figura, una tabla y un
párrafo por revisión. El protocolo completo, las tablas complementarias y la
discusión que sostienen ese resumen siguen a continuación.

\subsection{Protocolo comparado y capas del corpus}
\label{app:review-protocol-full}

La revisión académica se ejecuta como un embudo de cuatro capas, y cada capa
sostiene un tipo de afirmación distinto. La Tabla~\ref{tab:lit-corpus-layers}
lo hace explícito: el descubrimiento maximiza cobertura pero solo aporta
metadatos; el mapeo analítico permite hablar de tendencias e intersecciones;
el núcleo canónico admite comparar capacidades y arquitectura de información;
y solo la lectura cercana, sobre veinte trabajos, sostiene las afirmaciones
técnicas más fuertes.

% Tabla 1 del consolidado académico (`MROB_literature_review_reworked (2).pdf`,
% página 1). Reconstruida en LaTeX porque ese PDF no trae fuente: el `.tex`
% depositado junto a él corresponde a una versión anterior de la revisión.
\begin{table}[!htbp]
  \centering
  \caption{Capas del corpus y tipo de afirmación que puede sostener cada una.}
  \label{tab:lit-corpus-layers}
{\setstretch{1}\fontsize{7.6}{8.6}\selectfont\setlength{\tabcolsep}{3pt}\renewcommand{\arraystretch}{1.25}
\begin{tabularx}{\linewidth}{@{}J{2.5cm}K{1.0cm}J{5.1cm}Z@{}}
\toprule
\textbf{Capa} & \textbf{$n$} & \textbf{Uso principal} & \textbf{Límite de evidencia}\\
\midrule
\rowcolor{soft}
Descubrimiento & 3\,014 & Maximizar cobertura, enlazar versiones y alimentar el cribado. & Metadatos; no sostiene afirmaciones técnicas.\\
Mapeo analítico & 244 & Tendencias, cobertura temática, intersecciones y estructura bibliométrica. & Señales de título y resumen más extracción estructural; un cero significa «no identificado bajo esta codificación».\\
\rowcolor{soft}
Núcleo canónico & 59 & Comparación de capacidades, arquitectura de información y nivel de distribución. & Codificación explícita con \texttt{yes/partial/no/unclear}; no se infiere desde \emph{keywords}.\\
Lectura cercana & 20 & Antecedentes más próximos, contraejemplos y delimitación de la brecha. & Texto completo y evidencia localizada; base de las afirmaciones técnicas más fuertes.\\
\bottomrule
\end{tabularx}}
  \viusource{elaboración propia a partir del protocolo de revisión}
\end{table}


\begin{table}[!htbp]
  \centering
  \caption{Protocolo comparado de las tres revisiones que sostienen la brecha.}
  \label{tab:method-review-protocol}
  \fontsize{8.5}{10}\selectfont
  \renewcommand{\arraystretch}{1.04}
  \begin{tabularx}{\textwidth}{@{}>{\raggedright\arraybackslash}p{2.15cm} >{\raggedright\arraybackslash}p{3.05cm} >{\raggedright\arraybackslash}p{3.05cm} Y@{}}
    \toprule
    Dimensión & Académica & Industrial & Patentaria \\
    \midrule
    Unidad de análisis & documento arbitrado o preprint enlazado & sistema o despliegue con fuente oficial & publicación patentaria \\
    Recuperación & Crossref, OpenAlex, rutas OA y \emph{snowballing} & auditoría web de fabricantes e integradores & texto, CPC G05D 1/69--1/6987, solicitante y \emph{snowballing} \\
    Volumen & 3\,014 descubiertos, 244 cribados, 168 con texto adecuado, 59 canónicos & 10 casos principales y 4 de contexto & 169 filas y 167 registros analíticos \\
    Corte & sin corte duro; \emph{snowballing} cerrado en la etapa 5 & auditoría web de septiembre de 2026 & 11 de septiembre de 2026 \\
    Codificación & autoridad, información, líder y cierre, no solo control local & seis capacidades frente a la cadena operativa del TFM & una etiqueta temática y una macrocapa por publicación \\
    Qué certifica & mecanismos con formulación y garantía publicadas & capacidades disponibles en producto & desplazamiento de la actividad de protección \\
    Qué no certifica & prevalencia poblacional ni calidad experimental & cuota de mercado ni inexistencia de lo no documentado & número de invenciones ni familias jurídicas \\
    \bottomrule
  \end{tabularx}
  \viuownsource
\end{table}

«Distribuido» no se asigna por la presencia de control local, sino por la
combinación de autoridad de decisión, información disponible, existencia de
líder o servidor y responsable del cierre, de modo que un método con control
local y adjudicación centralizada se codifica como mixto. El corpus
patentario se analiza a nivel de publicación y no de invención ni de familia
jurídica, con el año de publicación como eje temporal porque el de prioridad
solo consta en 6 de los 167 registros. Cada publicación recibe además una
única etiqueta temática y una macrocapa ---decisión de misión, coordinación
cinemática o interacción física--- para impedir el doble conteo, conforme al
criterio de los manuales de análisis patentario
\parencite{wipoPatentLandscape,oecdPatentManual2009}, con las clases CPC
G05D~1/69, 1/692, 1/695, 1/696 y 1/698--1/6987 \parencite{usptoCpcG05D}.
El fichero \texttt{bibliography/references-v2.bib} conserva la ficha completa
de cada entrada catalogada durante las tres revisiones; el capítulo de
Referencias solo lista las que el argumento del texto cita explícitamente,
que es un subconjunto del corpus de 59, 14 y 167 declarado más arriba.

Las tres revisiones se proyectan sobre una misma cadena operativa. La
Figura~\ref{fig:method-review-sequence} fija sus seis etapas y anota bajo cada
una qué evidencia predomina hoy: las de decisión aparecen sobre todo en
investigación y patentes, la ejecución tiene precedentes comerciales
consolidados y la recomposición durante la misión solo acumula evidencia
parcial.

% Figura 3 del consolidado industrial, portada desde
% `thesis-reviews/TFM_estado_industrial_patentes_EDITABLE_v47.tex` (líneas
% 468-493). Sin cambios en el cuerpo TikZ.
\begin{figure}[!htbp]
  \centering
  \makebox[\textwidth][c]{\resizebox{\textwidth}{!}{%
\begin{tikzpicture}[x=0.985mm,y=1mm,font=\sffamily\fontsize{6.65}{7.15}\selectfont]
\path[use as bounding box] (0,0) rectangle (174,48);
\tikzset{stage/.style={draw=rule,line width=.45pt,minimum width=25mm,minimum height=13.5mm,align=left,fill=white,inner sep=0pt},arr/.style={-{Latex[length=1.7mm]},line width=.55pt,draw=onyx!45}}
% each stage is a white box with editable color rail
\foreach \x/\num/\title/\sub/\col in {14/1/{Reclutamiento}/{local por capacidad}/autumn,43/2/{Factibilidad}/{contacto / torsor}/autumn!82,72/3/{Aproximación}/{y formación}/mustard,101/4/{Transporte}/{carga compartida}/onyx!38,130/5/{Recomposición}/{ante fallo / cambio}/spicy!55,159/6/{Entrega}/{y disolución}/onyx!18}{
 \node[stage] (s\num) at (\x,31) {};
 \fill[\col] ({\x-12.5},24.25) rectangle ({\x-8.7},37.75);
 \node[font=\bfseries\fontsize{8.1}{8.5}\selectfont,text=white] at ({\x-10.6},31){\num};
 \node[anchor=west,font=\bfseries\fontsize{6.8}{7.2}\selectfont] at ({\x-6.8},32.6){\title};
 \node[anchor=west,text=muted] at ({\x-6.8},28.8){\sub};
}
\foreach \a/\b in {1/2,2/3,3/4,4/5,5/6}{\draw[arr] (s\a.east)--(s\b.west);}
% coverage labels
\node[anchor=north,text=muted,font=\fontsize{6.0}{6.4}\selectfont] at (14,22.7){I+D / patentes};
\node[anchor=north,text=muted,font=\fontsize{6.0}{6.4}\selectfont] at (43,22.7){I+D / patentes};
\node[anchor=north,text=muted,font=\fontsize{6.0}{6.4}\selectfont] at (72,22.7){literatura técnica};
\node[anchor=north,text=muted,font=\fontsize{6.0}{6.4}\selectfont] at (101,22.7){precedentes comerciales};
\node[anchor=north,text=muted,font=\fontsize{6.0}{6.4}\selectfont] at (130,22.7){evidencia parcial};
\node[anchor=north,text=muted,font=\fontsize{6.0}{6.4}\selectfont] at (159,22.7){operación comercial};
% three-layer band
\fill[autumn!24] (1,7.5) rectangle (61,11.5);\fill[mustard!42] (61,7.5) rectangle (113,11.5);\fill[onyx!18] (113,7.5) rectangle (173,11.5);
\node[anchor=north west,text=muted] at (1,7.1){decisión estratégica};\node[anchor=north west,text=muted] at (61,7.1){factibilidad y ejecución física};\node[anchor=north west,text=muted] at (113,7.1){continuidad de misión};
\node[anchor=west,text=muted,font=\fontsize{5.8}{6.2}\selectfont] at (1,2.2){Cadena evaluada: composición dinámica $\rightarrow$ verificación física $\rightarrow$ ejecución $\rightarrow$ recuperación.};
\end{tikzpicture}%
  }}
  \caption[Secuencia funcional y evidencia pública predominante.]{Secuencia
  funcional considerada en el trabajo y relación cualitativa con la evidencia
  pública predominante revisada. La banda inferior separa decisión estratégica,
  factibilidad y ejecución física, y continuidad de misión.}
  \label{fig:method-review-sequence}
  \viuownsource
\end{figure}


La codificación estructural del corpus académico es un primer pase
automatizado y no sustituye la lectura humana de ecuaciones, supuestos y
protocolos. La heterogeneidad de tareas, robots, contactos, comparadores y
métricas impide un metaanálisis cuantitativo. Ninguna de las tres revisiones
evaluó estabilidad, convergencia, optimalidad, seguridad o escalabilidad: esas
propiedades se contrastan, dentro de sus supuestos, en el capítulo de
resultados.

El mapa metodológico (Figura~\ref{fig:lit-methodological-map}) cubre 32 de
los 59 trabajos canónicos, con el criterio de selección del consolidado
original que le da fuente. Los 27 trabajos restantes no aparecen en el mapa,
aunque siguen contando en el corpus de 59 para la matriz de capacidades y la
auditoría adversarial de novedad; no se comprobó si su exclusión concentra
algún cuadrante en particular.

\subsection{Auditoría adversarial de novedad}
\label{app:review-redteam-full}

El primer frente resuelve la coalición lógica mediante emparejamientos uno a
muchos, subastas y juegos hedónicos
\parencite{choiBrunetHow2009CBBA,serviceAdams2011Coalition,chenSun2012Coalition,jangShinTsourdos2018GRAPE,duttaAsaithambi2019Bipartite,dutta2021hedonic,diehlAdams2026Grapes};
el segundo resuelve la ejecución física con formulación, control y validación
experimental
\parencite{alonsomora2017transport,shibata2023event,ebel2024cooperative,fukao2025swarm,tian2025irregular,muhammed2026decentralizedMpc}.
La brecha no puede formularse como «falta control distribuido de carga», que
sería falso, sino como una interfaz que ninguno de los dos frentes atraviesa.

Tres lecturas condicionan el diseño del resto de la memoria. La atomicidad
exige que el cierre entero pertenezca al método, porque los juegos
poblacionales producen masa continua y un redondeo posterior puede romper la
propiedad demostrada en la Ecuación~\eqref{eq:tf-population-dynamics}. El paso
de capacidad a \emph{wrench} impide sustituir
$W_k^{\mathrm{req}} \in \mathcal W_C(q)$ por la suma escalar
$\sum_i c_i \ge d_k$, según la Ecuación~\eqref{eq:tf-grasp-map}. La escala
obliga a separar robots simulados, robots físicos, robots por coalición y
cargas simultáneas. A esto se añade la modalidad de interacción: carga
soportada, acoplamiento rígido, empuje y prensión imponen conjuntos de
esfuerzos admisibles distintos y no deben agruparse (panel inferior de la
Figura~\ref{fig:lit-capability-matrix}).

% Tabla 2 del consolidado académico (PDF vigente, página 4): antecedentes más
% próximos frente a las ocho capacidades que definen la combinación objetivo.
% Reconstruida en LaTeX; el `.tex` depositado corresponde a una versión previa
% con otra tabla de red-team.
\begin{table}[!htbp]
  \centering
  \caption{Antecedentes más próximos frente a las capacidades que definen la combinación objetivo.}
  \label{tab:lit-prior-art}
{\setstretch{1}\fontsize{7.6}{8.6}\selectfont\setlength{\tabcolsep}{3.4pt}\renewcommand{\arraystretch}{1.22}
\rowcolors{2}{white}{soft}
\begin{tabularx}{\linewidth}{@{}J{2.6cm}*{8}{K{0.62cm}}Z@{}}
\toprule
\textbf{Trabajo} & \textbf{H} & \textbf{V} & \textbf{L} & \textbf{S} & \textbf{W} & \textbf{D} & \textbf{R} & \textbf{X} & \textbf{Restricción decisiva}\\
\midrule
Verma 2019 \parencite*{verma2019replacement} & \capno & \cappart & \capno & \capyes & \cappart & \capno & \capyes & \capyes & planificación global; ruta e hitos conocidos\\
Sahu 2026 \parencite*{sahuKumar2026FaultManagement} & \capno & \cappart & \capno & \capyes & \cappart & \cappart & \capyes & \capyes & arquitectura maestro--esclavo\\
Shibata 2023 \parencite*{shibata2023event} & \capno & \cappart & \capno & \capyes & \capyes & \capyes & \cappart & \cappart & equipo variable; no selección desde una flota heterogénea\\
Ebel 2024 \parencite*{ebel2024cooperative} & \capno & \capno & \capno & \capyes & \capyes & \capyes & \capno & \capno & equipo de transporte fijado\\
Bezerra 2025 \parencite*{bezerra2025dynamicCoalition} & \capyes & \capyes & \capyes & \capno & \capno & \capyes & \cappart & \capno & reclutamiento lógico; sin carga compartida física\\
GRAPE-S 2026 \parencite*{diehlAdams2026Grapes} & \capyes & \capyes & \cappart & \capno & \capno & \capyes & \capno & \capno & formación de servicios; no transporte de carga compartida\\
\bottomrule
\end{tabularx}}
\par\vspace{1.5pt}
{\setstretch{1}\fontsize{7.0}{7.8}\selectfont\color{muted}
H = heterogeneidad; V = tamaño de coalición variable; L = reclutamiento local; S = carga compartida; W = certificado físico; D = ejecución o decisión distribuida; R = recuperación; X = sustitución de miembro. \capyes\ explícito, \cappart\ parcial o condicionado, \capno\ no documentado en el alcance revisado.\par}
  \viusource{codificación del núcleo canónico y evidencia de texto completo}
\end{table}


Una brecha solo es defendible si sobrevive al intento de refutarla, y la
Tabla~\ref{tab:lit-prior-art} recoge esa auditoría adversarial. El resultado
obliga a estrechar el enunciado: el reemplazo durante el transporte ya existe.
Verma et al. lo implementan con ruta y puntos de relevo conocidos y
planificación global \parencite{verma2019replacement}; Sahu y Kumar lo
integran con detección de fallo bajo arquitectura maestro--esclavo
\parencite{sahuKumar2026FaultManagement}; y Shibata et al. admiten que agentes
se retiren o incorporen durante el control distribuido, sin resolver la
selección del sustituto dentro de una flota heterogénea
\parencite{shibata2023event}. Ebel et al. resuelven el transporte cooperativo
con un equipo fijado de antemano, sin mecanismo de reclutamiento ni de
sustitución \parencite{ebel2024cooperative}. En el extremo opuesto, Bezerra y
GRAPE-S cubren reclutamiento o formación de coaliciones sin cerrar la
interacción física de una carga compartida.

El enunciado de brecha reclama solo que la composición de las tres piezas bajo
información local carece de precedente documentado en el corpus: no el primer
método distribuido de transporte cooperativo, ni el primer mecanismo de
reemplazo, ni superioridad frente a ningún trabajo de la
Figura~\ref{fig:lit-capability-matrix}.

\subsection{Panorama industrial y marco normativo}
\label{app:review-industrial-full}

La lectura relevante del mapa industrial es que dos líneas funcionales
coexisten con precedentes comerciales propios. En el cuadrante inferior
derecho, la cooperación física para carga pesada está resuelta con
configuración de ingeniería y escasa autonomía: KUKA omniMove mueve secciones
de fuselaje de hasta 100~toneladas en Airbus Stade con composición dedicada
\parencite{kukaOmnimoveAirbus2016} y los transportadores modulares SPMT
resuelven \emph{project cargo} mediante configuración previa
\parencite{tiiScheuerleSpmt2023}. En el superior izquierdo, la autonomía de
flota está igualmente resuelta pero sin carga compartida: AGILOX opera
27~AMR en BMW Regensburg con coordinación entre pares y sin control central
\parencite{agiloxXswarmBmw2023}. El cuadrante superior derecho, donde ambas
capacidades convergen, permanece escasamente poblado.

Beacon Dual-AMR es el precedente con mayor solapamiento funcional, porque
documenta emparejamiento automático y relevo de autoridad ante fallo, pero
opera con dos robots bajo despacho central y no sustituye a un miembro por
otro de la flota \parencite{beaconDualAmr2026}. Dos heterogeneidades distintas
conviven en este panorama. La de flota, entendida como convivencia de modelos
distintos bajo un mismo gestor, aparece con frecuencia en las plataformas de
contexto. La que se da dentro de una misma coalición física, donde robots con
capacidades y geometrías diferentes deben cerrar conjuntamente restricciones
de carga, contacto y \emph{wrench}, tiene cobertura pública mucho menor y es
la que el problema focal aborda.

% Tabla 2 del consolidado industrial, portada desde
% `thesis-reviews/TFM_estado_industrial_patentes_EDITABLE_v47.tex` (líneas
% 498-509). Las marcas \casecite se sustituyen por citas biblatex reales y
% los tipos de columna L{}/Y pasan a J{}/Z.
\begin{table}[!htbp]
  \centering
  \caption{Normas e interfaces relevantes para la integración industrial.}
  \label{tab:tf-standards}
{\setstretch{1}\fontsize{6.95}{7.65}\selectfont\setlength{\tabcolsep}{2.2pt}\renewcommand{\arraystretch}{1.16}
\begin{tabularx}{\linewidth}{@{}J{3.15cm}Z@{}}
\toprule
\textbf{Documento} & \textbf{Relación con el problema}\\
\midrule
\rowcolor{soft} VDA 5050 v3.0.0 (2026) & Interfaz de órdenes y estado entre control de flota y robots móviles. La v3.0 incorpora zonas de libre navegación y compartición de rutas (\textit{path sharing}); favorece interoperabilidad operativa, pero no define la constitución ni la verificación mecánica de una coalición de carga compartida \parencite{vda5050v3}.\\
ISO 21423 (2026, en publicación) & Protocolos de comunicación para interoperabilidad entre AMR multivendor, gestores de flota y recursos empresariales. Consultado el 11-09-2026: figura en etapa 60.00, ``International Standard under publication'', y excluye requisitos de seguridad \parencite{iso21423}.\\
\rowcolor{soft} ISO 3691-4:2023 & Requisitos de seguridad y medios de verificación para vehículos industriales sin conductor y sus sistemas; incluye AGV y AMR como ejemplos \parencite{iso36914}.\\
ANSI/A3 R15.08-3:2026 & Uso seguro de aplicaciones IMR: evaluación de riesgos, gestión de cambios y mantenimiento del riesgo aceptable durante el ciclo de vida \parencite{ansiA3R1508part3}.\\
\bottomrule
\end{tabularx}}
  \viusource{elaboración propia a partir de los documentos citados}
\end{table}


La Tabla~\ref{tab:tf-standards} delimita el marco normativo de cualquier
realización industrial de esta arquitectura. Ni VDA 5050 ni ISO 21423 definen
la constitución de una coalición física o su certificación mecánica, e ISO
3691-4 y ANSI/A3 R15.08-3 fijan el marco de seguridad sin cubrir ese contrato.
Queda por tanto fuera de la normativa vigente y debe resolverlo el sistema,
con una consecuencia de diseño directa: el certificado de \emph{wrench} no
puede delegarse en una capa de interoperabilidad y debe producirse dentro del
propio mecanismo de reclutamiento.

\subsection{Discusión ampliada de cobertura, tendencia y patentes}
\label{app:review-results-full}

Esta subsección conserva, sin abreviar, la lectura panel a panel de las
figuras de cobertura, tendencia metodológica y actividad patentaria que la
Sección~\ref{subsec:results-review} presenta de forma condensada.

El panel (a) de la Figura~\ref{fig:lit-coverage-rubric} sostiene el planteamiento: el
corpus cubre con solvencia las capacidades por separado ---veintitrés trabajos abordan
transporte compartido y dieciocho el reclutamiento--- y la cobertura se desploma al
exigir intersecciones, con solo dos trabajos que combinan carga compartida y reemplazo y
ninguno que reúna certificación física, ejecución distribuida y reemplazo. El panel (b)
explica cómo se codificó la dependencia global: la rúbrica se aplica trabajo a trabajo y
etapa a etapa, sin promediar categorías ordinales, y es la que permite afirmar que la
autoridad residual se desplaza hacia lo global precisamente en certificación y
recuperación.

La Figura~\ref{fig:lit-stage-coverage} añade el estrechamiento decisivo. El panel (a)
muestra que la cobertura se concentra en SP1, con 115 señales frente a 20 de SP3 y 15 de
SP2. El panel (b) muestra el segundo estrechamiento: 103 documentos se codifican
únicamente en SP1, diez enlazan SP1--SP3, dos enlazan SP1--SP2 y ninguno cubre las tres
etapas. El panel (c) explica por qué los comparadores deben elegirse con cuidado: hay 99
señales de simulación, 25 de contexto o caso industrial, 12 de experimento físico y
siete de análisis formal. Esa proporción de validación física es la razón de que la
memoria mantenga separados los resultados de simulación y los del demostrador.

La Figura~\ref{fig:lit-method-change} descompone el cambio metodológico en cuatro
periodos y muestra recombinación, no sustitución. Consenso y distribución pasan del
20,0~\% al 24,7~\% y subasta o mercado baja del 22,9~\% al 13,5~\%, pero ambas familias
siguen presentes; en paralelo, el aprendizaje sube 12,8~puntos, la metaheurística
evolutiva 11,2 y la seguridad reactiva o CBF 6,7. El patrón dominante es de
recombinación: selección, coordinación, control y seguridad ganan diversidad sin
converger hacia una única familia metodológica.

La Figura~\ref{fig:lit-period-shift} cuantifica ese mismo desplazamiento sobre las
cinco categorías verificables del corte V3. La cuota de contexto general crece 17,7
puntos entre el periodo temprano y 2022--2026, y la de recuperación o resiliencia cae
9,7; coalición y transporte se mantienen casi estables, con cambios menores a un
punto. La lectura es la misma que en el panel de tendencias: el corpus no converge
hacia una única familia, sino que redistribuye peso entre las que ya tenía.

La estructura bibliométrica de la Figura~\ref{fig:lit-bibliometric} confirma esa
lectura desde otro ángulo. Los núcleos de coautoría fuerte son pequeños y aislados
---como máximo cuatro autores---, de modo que ninguna agenda de investigación domina
el corpus recuperado. Las comunidades de método tienen modularidad baja
($Q=$\LitMethodModularity{}), señal de que consenso, subasta, aprendizaje y
metaheurística se combinan con frecuencia en un mismo trabajo en vez de segregarse en
escuelas cerradas.

El corpus patentario responde a una pregunta distinta: hacia dónde se desplaza el
interés de protección industrial. Con la media móvil de tres años calculada al cierre
de cada periodo (2022 y 2025 respectivamente), la macrocapa de decisión
de misión pasa del 37,9~\% al 58,8~\% del corpus, un aumento de 20,9~puntos porcentuales,
mientras la coordinación cinemática baja del 30,3~\% al 17,6~\% y la interacción física
del 31,8~\% al 23,5~\%. El resultado no depende de la ampliación dirigida por
solicitante: excluyéndola el cambio va del 32,7~\% al 54,7~\%, es decir 22,0~puntos, y la
prueba de sensibilidad que colapsa los cuatro grupos de títulos muy similares lo deja en
20,0~puntos. Los paneles (b) y (c) añaden dos matices para no sobreinterpretarlo: China
concentra la mayoría de titulares identificados y el 14,3~\% reciente de Estados Unidos
procede íntegramente de la ampliación dirigida, de modo que no admite lectura como cuota
mundial; por dominios, la intralogística concentra el 69~\% de su actividad en decisión
de misión y el sector aeroespacial el 88~\% en interacción física.
